\clearpage
\FrontTitle{ABSTRAK}
\begin{center}
{\bfseries\normalsize \JudulLaporan\par}
\vspace{0.1cm}
{\bfseries\normalsize \TempatKP\par}
\vspace{0.2cm}
{Oleh\par
\bfseries\normalsize \NamaMahasiswa\par
\bfseries\normalsize NIM: \NIM\par}
\end{center}
\vspace{0.4cm}

Kerja Praktik dilaksanakan di PT Paragon Technology and Innovation pada Direktorat \textit{Information Technology} (tim FATL) selama 10 minggu sejak 2 Juni 2026 hingga 11 Agustus 2026. Penulis berperan sebagai \textit{Software Engineer} yang berfokus pada pengembangan sistem Parapay, sebuah aplikasi manajemen dan penagihan piutang usaha (\textit{Account Receivable}) yang digunakan oleh kolektor FAR (\textit{Field Account Receivable}). Pengelolaan penagihan piutang pada perusahaan berskala besar menghadapi tantangan kompleksitas dalam pencatatan alokasi pembayaran multi-transaksi, sinkronisasi alur revisi pembayaran kas maupun non-kas, penanganan bukti bayar secara luring, serta beban data transaksi dalam jumlah besar. Untuk mengatasi permasalahan tersebut, dilakukan pengembangan beberapa fitur utama: implementasi alur simpan draf (\textit{Save as Draft}), modul Potong Tagihan, perombakan arsitektur alur revisi dari Revision v1 menjadi Revision v2 yang modular dengan pembaruan status otomatis (\textit{state invalidation}), integrasi penyimpanan bukti transfer berbasis \textit{IndexedDB} untuk persistensi berkas secara luring, serta optimasi antarmuka melalui \textit{server-side search} dan \textit{infinite scroll}. Pengembangan dilaksanakan menggunakan pendekatan \textit{Agile} dengan siklus mingguan melalui \textit{daily standup} dan pengujian fungsional terintegrasi. Hasil Kerja Praktik ini berupa terwujudnya alur revisi pembayaran yang transparan dan andal, berkurangnya kegagalan unggah bukti bayar, serta peningkatan responsivitas antarmuka Parapay dalam mengelola ribuan data faktur pelanggan.

\vspace{0.3cm}
\noindent\textbf{Kata kunci}: \textit{Account Receivable}, \textit{Revision Flow}, \textit{IndexedDB}, \textit{State Invalidation}, \textit{Frontend Engineering}, Parapay.
